%% maestro2tex -- generated figure; do not edit by hand.
\providecommand{\MaestroRoot}{include/analog/}
\begin{figure}[htbp]
  \centering
  {\pgfplotsset{compat=1.18}%
  \begin{tikzpicture}
    \begin{axis}[
      width=0.80\linewidth, height=6.2cm,
      xlabel={Commanded output voltage $V_{\mathrm{MAP}}$~[\si{\volt}]}, ylabel={Compliance span~[\si{\micro\ampere}]},
      grid=both,
      major grid style={line width=0.2pt, draw=black!14},
      minor grid style={line width=0.2pt, draw=black!7},
      minor tick num=1,
      axis line style={line width=0.4pt, draw=black!45},
      tick style={line width=0.4pt, draw=black!45},
      tick align=outside,
      label style={font=\small}, tick label style={font=\small},
      enlarge x limits=0.02, enlarge y limits=0.10,
      every axis plot/.append style={line join=round, no marks},
      legend style={font=\footnotesize, draw=none, fill=none,
                    at={(0.5,1.02)}, anchor=south, legend columns=-1,
                    /tikz/every even column/.append style={column sep=9pt}},
      legend cell align=left,
      scaled y ticks=false, scaled x ticks=false,
      y tick label style={/pgf/number format/fixed,
                          /pgf/number format/precision=1},
      ymin=0, ymax=16, enlarge y limits=false, restrict y to domain=0:60
    ]
      \addplot[mtxA, line width=1.1pt, solid] table {\MaestroRoot data/ctrlamp_drivespan_00.dat};
      \addlegendentry{tt}
      \addplot[mtxB, line width=1.1pt, dashed] table {\MaestroRoot data/ctrlamp_drivespan_01.dat};
      \addlegendentry{ff}
      \addplot[mtxC, line width=1.1pt, dotted] table {\MaestroRoot data/ctrlamp_drivespan_02.dat};
      \addlegendentry{fs}
      \addplot[mtxD, line width=1.1pt, dashdotted] table {\MaestroRoot data/ctrlamp_drivespan_03.dat};
      \addlegendentry{sf}
      \addplot[mtxE, line width=1.1pt, densely dashed] table {\MaestroRoot data/ctrlamp_drivespan_04.dat};
      \addlegendentry{ss}
    \end{axis}
  \end{tikzpicture}}
  \caption{Width of the \(\pm\SI{50}{\milli\volt}\) compliance envelope against commanded output voltage, over all five process corners. The shape is the same in every corner --- pinched at both ends of the output range and widest around \SI{1.8}{\volt} --- so the deficit against the \(\pm\SI{33}{\micro\ampere}\) target is structural, not a corner artefact. The vertical scale is clipped at \SI{16}{\micro\ampere}: outside roughly \SIrange{0.3}{1.9}{\volt} the buffer no longer tracks its input at zero load, and the apparent span there rises past \SI{30}{\micro\ampere} because it is measuring the current needed to force the output onto the commanded value, not compliance. Note also that the resistor and MIM rows are pinned at typical in all five corners (see the note below), so this is the transistor-corner spread only and understates the true one.}
  \label{fig:ctrlamp-drivespan}
\end{figure}
